\section{Theory of Constrained Optimization}
The formulation of constrained problems in this book,
\[
    \underset{x\in \mathbb{R}^n}{\min} f(x) \qquad \text{subject to} \quad 
    \begin{cases}
        c_i(x) &= 0, \quad i\in\mathcal{E},\\
        c_i(x) &\geq 0, \quad i\in\mathcal{I}
    \end{cases}
\]
is standard in the literature, except for the sign chosen of the inequality constraint functions ($i\in\mathcal{I}$) which may vary from author to author. 
By defining the feasible set as
\[
    \Omega \equiv \left\{ x \in \mathbb{R}^n | c_i(x) = 0, \forall i\in\mathcal{E}, c_i(x) \geq 0, \forall i \in \mathcal{I} \right\}
\]
the authors present a very elegant formulation of the constrained problem as
\[
    \underset{x\in\Omega}{\min}f(x).
\]
% As a fun exercise, show that the tangent cone is a vector space over R.
% Also, it should be possible to show it is a smooth manifold. Does it need several charts?
% x1^2 + x2^2 < 2 needs two charts. Probably just need to define a transition function, then done!
% -------------------------------------
% #####################################
% ------------ Cone Stuff -------------
% #####################################
% -------------------------------------
\subsection{Solutions}
The exercises chosen for this chapter are 12.1, 12.2, 12.4, 12.4, 12.11 (not \textbf{b)}) and 12.22. 